% Part: turing-machines
% Chapter: machines-computations
% Section: representing-tms

\documentclass[../../../include/open-logic-section]{subfiles}

\begin{document}

\olfileid{tur}{mac}{rep}
\olsection{د ټيورينګ ماشينونو څرګندونه}

\begin{explain}
ټيورينګ ماشينونه په ليديز ډول د \emph{حالت ډياګرامونو} په وسيله
ښودل کېدے شي۔ دا ډياګرامونه د حالت له هغو خانو جوړ دي چې غشي ئې سره
نښلوي۔ لکه چې تمه کېږي، د حالت هره خانه د ماشين د يو حالت استازيتوب
کوي۔ هر غشی د هغې لارښوونې استازيتوب کوي چې له هماغه حالت څخه عملي
کېدے شي، او د لارښوونې جزئيات د اړوند غشي دپاسه يا لاندې ليکل شوي
وي۔ لاندې ماشين وګورئ، چې يوازې دوه دننني حالتونه، $q_0$ او~$q_1$،
او يوه لارښوونه لري:
\[
\begin{tikzpicture}[->,>=stealth',shorten >=1pt,auto,node distance=2.8cm,
                    semithick]
  \tikzstyle{every state}=[fill=none,draw=black,text=black]

  \node[initial,state]   (A)              {$q_0$};
  \node[state]   (B) [right of=A] {$q_1$};

  \path (A) edge  node {\TMtrans{\TMblank}{\TMstroke}{\TMright}} (B);
\end{tikzpicture}
\]
په ياد ولرئ چې ټيورينګ ماشين يو لوست-ليک سر او يوه پټه لري چې ننوت
پرې ليکل شوے وي۔ لارښوونه داسې لوستل کېږي: \emph{که په حالت~$q_0$
کښې~$\TMblank$ لولئ، يو~$\TMstroke$ وليکئ، ښي لوري ته لاړ شئ، او
حالت~$q_1$ ته واوړئ}۔ دا له دې سره معادل ده چې د حالت بدلون تابع
$\tuple{q_0, \TMblank}$ په~$\tuple{q_1, \TMstroke,
\TMright}$ واړوي۔
\end{explain}

\begin{ex}
\emph{جفت ماشين}: لاندې ټيورينګ ماشين هله او يوازې هله درېږي چې پر
پټه د~$\TMstroke$ ګانو شمېر جفت وي (په دې فرض چې د پټې ټول
$\TMstroke$ ګانې له لومړي~$\TMblank$ څخه مخکښې راځي)۔
\[
\begin{tikzpicture}[->,>=stealth',shorten >=1pt,auto,node distance=2.8cm,
                    semithick]
  \tikzstyle{every state}=[fill=none,draw=black,text=black]

  \node[initial,state]         (A)              {$q_0$};
  \node[state]         (B) [right of=A] {$q_1$};

  \path (A) edge [bend left] node {\TMtrans{\TMstroke}{\TMstroke}{\TMright}} (B)
        (B) edge [loop above] node {\TMtrans{\TMblank}{\TMblank}{\TMright}} (B)
            edge [bend left] node {\TMtrans{\TMstroke}{\TMstroke}{\TMright}} (A);
\end{tikzpicture}
\]

د حالت ډياګرام له لاندې حالت بدلون تابع سره سمون لري:
\begin{align*}
\delta(q_0, \TMstroke) & = \tuple{q_1, \TMstroke, \TMright},\\
\delta(q_1, \TMstroke) & = \tuple{q_0, \TMstroke, \TMright},\\
\delta(q_1, \TMblank)  & = \tuple{q_1, \TMblank, \TMright}
\end{align*}
\end{ex}

\begin{explain}
پورته ماشين يوازې هغه وخت درېږي چې ننوت د کرښو جفت شمېر وي۔ که نه،
ماشين (په نظري ډول) تر نامعلومې مودې کار ته دوام ورکوي۔ د هر ماشين
او ننوت دپاره د ماشين \emph{تشکيلونه} پرله‌پسې څارلے شو، څو وت ئې
وټاکو۔ د تشکيلونو صوري تعريف به وروسته ورکړو۔ اوس د تشکيلونو په اړه
په شهودي ډول د ډياګرامونو د يوې لړۍ په توګه فکر کولے شو، چې د کار په
بهير کښې په هره شېبه د ماشين حالت ښيي۔ تشکيلونه د پټې منځپانګه، د
ماشين حالت او د لوست-ليک سر ځاے ښيي۔

راځئ د جفت ماشين تشکيلونه هغه وخت پرله‌پسې وڅارو چې د څلورو
$\TMstroke$ ګانو پر ننوت پيل شي۔ په دې صورت کښې تمه لرو چې ماشين
ودرېږي۔ بيا به ماشين د درېيو~$\TMstroke$ ګانو پر ننوت وچلوو، چې
هلته به ماشين تل روان وي۔

ماشين په حالت~$q_0$ کښې، د تر ټولو کيڼ~$\TMstroke$ په لوستلو پيل
کوي۔ د ماشين لومړنی حالت داسې ښودلے شو:
\[
\TMendtape \TMstroke_0 \TMstroke \TMstroke \TMstroke \TMblank \ldots
\]
پورته تشکيل روښانه دے۔ لکه چې ليدل کېږي، ماشين په صفرم حالت کښې د
تر ټولو کيڼ~$\TMstroke$ په لوستلو پيل کوي۔ دا پر لومړي~$\TMstroke$
د حالت د نوم په فرعي نښه ښودل کېږي۔
\textbf{د ژباړې د نسخې سمون (OLCMP-040):} سرچينې په دې ځاے کښې د
ماشين حالت په نثر کښې لومړی حالت بللے و، خو ډياګرام او فرعي نښه دواړه
صفرم حالت ټاکي؛ دلته هماغه سم حالت راوړل شوے دے۔
په دې ځاے کښې د عملي کېدو وړ لارښوونه
$\delta(q_0, \TMstroke) = \tuple{q_1, \TMstroke, \TMright}$ ده، نو
ماشين پر پټه ښي لوري ته ځي او حالت~$q_1$ ته اوړي۔
\[
\TMendtape \TMstroke \TMstroke_1 \TMstroke \TMstroke \TMblank \ldots
\]
څرنګه چې ماشين اوس په حالت~$q_1$ کښې يو~$\TMstroke$ لولي، بايد
لارښوونه~$\delta(q_1, \TMstroke) = \tuple{q_0,
  \TMstroke, \TMright}$ ``تعقيب'' کړو۔ په نتيجه کښې دا تشکيل لاس ته
راځي:
\[
\TMendtape \TMstroke \TMstroke \TMstroke_0 \TMstroke \TMblank \ldots
\]
د ماشين د دوام سره قواعد يو ځل بيا په هماغه ترتيب عملي کېږي، او
لاندې دوه تشکيلونه لاس ته راځي:
\[
\TMendtape \TMstroke \TMstroke \TMstroke \TMstroke_1 \TMblank \ldots
\]
\[
\TMendtape \TMstroke \TMstroke \TMstroke \TMstroke \TMblank_0 \ldots
\]
ماشين اوس په حالت~$q_0$ کښې يو~$\TMblank$ لولي۔ د حالت بدلون د
ډياګرام له مخې په اسانه وينو چې د عملي کولو هېڅ لارښوونه نشته، نو
ماشين درېدلے دے۔ معنا دا چې ننوت منل شوے دے۔

اوس فرض کړئ ماشين د درېيو~$\TMstroke$ ګانو پر ننوت پيلوو۔ لومړني
څو تشکيلونه ورته دي، ځکه هماغه لارښوونې عملي کېږي او يوازې د پټې په
ننوت کښې لږ توپير شته:
\[
\TMendtape \TMstroke_0 \TMstroke \TMstroke \TMblank \ldots
\]
\[
\TMendtape \TMstroke \TMstroke_1 \TMstroke \TMblank \ldots
\]
\[
\TMendtape \TMstroke \TMstroke \TMstroke_0 \TMblank \ldots
\]
\[
\TMendtape \TMstroke \TMstroke \TMstroke \TMblank_1 \ldots
\]
ماشين اوس له ټولو~$\TMstroke$ ګانو څخه تېر شوے، او په حالت~$q_1$
کښې يو~$\TMblank$ لولي۔ لکه چې په ډياګرام کښې ښودل شوې، د
$\delta(q_1, \TMblank) =\tuple{q_1,
\TMblank, \TMright}$ په بڼه يوه لارښوونه شته۔ څرنګه چې پټه ښي لوري
ته تر نامعلومې مودې پورې له~$\TMblank$ څخه ډکه ده، ماشين به دا
لارښوونه \emph{تل} عملي کوي، په حالت~$q_1$ کښې به پاتې کېږي او لا
پسې به ښي لوري ته ځي۔ ماشين به هېڅکله ونۀ درېږي، او ننوت نۀ مني۔
\end{explain}

\begin{explain}
دا خبره مهمه ده چې ټول ماشينونه نۀ درېږي۔ که درېدل په دې معنا وي چې
د ماشين د عملي کولو لارښوونې خلاصې شي، نو داسې ماشين جوړولے شو چې
هېڅکله نۀ درېږي: بس دا باوري کړو چې په هر حالت کښې د هرې نښې دپاره
يو وتونکے غشی شته۔ په حالت~$q_0$ کښې د~$\TMblank$ د لوستلو يوه
لارښوونه په ور زياتولو جفت ماشين داسې بدلولے شو چې تل روان وي۔
\end{explain}

\begin{ex}
\[
\begin{tikzpicture}[->,>=stealth',shorten >=1pt,auto,node distance=2.8cm,
                    semithick]
  \tikzstyle{every state}=[fill=none,draw=black,text=black]

  \node[initial,state] (A)              {$q_0$};
  \node[state]         (B) [right of=A] {$q_1$};

  \path
  (A) edge [bend left] node {\TMtrans{\TMstroke}{\TMstroke}{\TMright}} (B)
      edge [loop above] node {\TMtrans{\TMblank}{\TMblank}{\TMright}} (A)
  (B) edge [loop above] node {\TMtrans{\TMblank}{\TMblank}{\TMright}} (B)
      edge [bend left] node {\TMtrans{\TMstroke}{\TMstroke}{\TMright}} (A);
\end{tikzpicture}
\]
\end{ex}

\begin{explain}
د ماشين جدولونه د ټيورينګ ماشينونو د څرګندولو بله لار ده۔ د ماشين
جدولونه د پټې الفبا پر~$x$-محور، او د ماشين د حالتونو سټ پر~$y$-محور
ښيي۔ د جدول دننه، د هر حالت او نښې په ګډ ځاے کښې د لارښوونې پاتې
برخه---نوې نښه، نوی حالت او د خوځښت لوری---ليکل کېږي۔ د ماشين جدولونه
دا اسانه کوي چې ماشين په کوم حالت او د کومې نښې دپاره درېږي۔ په جدول
کښې هر تش ځاے د ماشين د درېدو يو ممکن ټکے دے۔ د حالت د ډياګرامونو او
د لارښوونو د سټونو پر خلاف، چې د ماشين د درېدو ځايونه پکښې تل سملاسي
څرګند نۀ وي، د ماشين په جدول کښې د تشو ځايونو په موندلو هر د درېدو
ټکے ژر پېژندل کېږي۔
\textbf{د ژباړې د نسخې سمون (OLCMP-041):} سرچينې د جدول د حجرې
توکي د نوي حالت او نوې نښې په سرچپه ترتيب ياد کړي وو، حال دا چې
ښودل شوې حجرې لومړے نوې نښه او بيا نوی حالت ليکي؛ د هر تش ځاے په
اړه ورپسې جمله کښې ورکه تړونکې کلمه هم ور زياته شوه۔
\end{explain}

\begin{ex}
د جفت ماشين جدول دا دے:
\[
\centering
\begin{tabular}{lllll}
\cline{2-4}
\multicolumn{1}{l|}{}      & \multicolumn{1}{c|}{$\TMblank$}
& \multicolumn{1}{c|}{$\TMstroke$}     & \multicolumn{1}{c|}{$\TMendtape$}   &  \\ \cline{1-4}
\multicolumn{1}{|l|}{$q_0$} & \multicolumn{1}{l|}{}
& \multicolumn{1}{l|}{\TMtrans{\TMstroke}{q_1}{\TMright}} &
\multicolumn{1}{l|}{\phantom{\TMtrans{\TMstroke}{q_1}{\TMright}}} &  \\ \cline{1-4}
\multicolumn{1}{|l|}{$q_1$} & \multicolumn{1}{l|}{\TMtrans{\TMblank}{q_1}{\TMright}}
& \multicolumn{1}{l|}{\TMtrans{\TMstroke}{q_0}{\TMright}} &
\multicolumn{1}{l|}{\phantom{\TMtrans{\TMstroke}{q_1}{\TMright}}} &  \\ \cline{1-4}
\end{tabular}
\]
لکه چې وينو، ماشين هغه وخت درېږي چې په حالت~$q_0$ کښې يو~$\TMblank$ لولي۔
\end{ex}

\begin{explain}
تر اوسه مو يوازې هغه ماشينونه کتلي چې ننوت لولي او مني ئې۔ خو
ټيورينګ ماشينونه د لوستلو او ليکلو دواړو وړتيا لري۔ د داسې ماشين يوه
بېلګه (که څه هم ډېرې، ډېرې بېلګې شته) \emph{دوه‌چنده کوونکے} دے۔
دوه‌چنده کوونکے چې پر پټه د~$n$ شمېر~$\TMstroke$ ګانو له يوې ډلې سره
پيل شي، د~$2n$ شمېر~$\TMstroke$ ګانو يوه ډله د وت په توګه ورکوي۔
\end{explain}

\begin{ex}
  \ollabel{ex:doubler} مخکښې له دې چې دوه‌چنده کوونکے ماشين جوړ کړو،
  د مسئلې د حل دپاره يوه \emph{تګلاره} جوړول مهم دي۔ څرنګه چې ماشين
  (لکه چې موږ تعريف کړے) دا نۀ شي ياد ساتلے چې څو~$\TMstroke$ ګانې
  ئې لوستې دي، بايد د پټې د ټولو~$\TMstroke$ ګانو د شمېر ساتلو يوه
  لار پيدا کړو۔ يوه داسې لار دا ده چې وت له ننوت څخه په يو~$\TMblank$
  جلا کړو۔ بيا ماشين له ننوت څخه لومړے~$\TMstroke$ پاکولے، د ننوت پر
  پاتې برخې تېرېدلے، يو~$\TMblank$ پرېښودلے، او دوه نوي~$\TMstroke$
  ګانې ليکلے شي۔ بيا به ماشين بېرته لاړ شي، په ننوت کښې به دويم
  $\TMstroke$ پيدا کړي، او هغه به هم دوه‌چنده کړي۔ د ننوت د هر يو
  $\TMstroke$ په بدل کښې به د وت دوه~$\TMstroke$ ګانې وليکي۔ د
  ماشين د تګ په ترڅ کښې د ننوت په پاکولو باوري کولے شو چې هېڅ
  $\TMstroke$ پاتې نۀ شي او نۀ دوه ځله دوه‌چنده شي۔ کله چې ټول ننوت
  پاک شي، پر پټه به~$2n$ شمېر~$\TMstroke$ ګانې پاتې وي۔ د لاس ته
  راغلي ټيورينګ ماشين د حالت ډياګرام په~\olref{fig:doubler} کښې
  ښودل شوے دے۔
  \begin{figure}
\[
\begin{tikzpicture}[->,>=stealth',shorten >=1pt,auto,node distance=2.8cm,
                    semithick]
  \tikzstyle{every state}=[fill=none,draw=black,text=black]

  \node[initial,state] (1)              {$q_0$};
  \node[state]         (2) [right of=1] {$q_1$};
  \node[state]         (3) [right of=2] {$q_2$};
  \node[state]         (4) [below of=3] {$q_3$};
  \node[state]         (5) [left of=4]  {$q_4$};
  \node[state]         (6) [left of=5]  {$q_5$};

  \path (1) edge node {\TMtrans{\TMstroke}{\TMblank}{\TMright}} (2)
    (2) edge [loop above] node {\TMtrans{\TMstroke}{\TMstroke}{\TMright}} (2)
      edge node {\TMtrans{\TMblank}{\TMblank}{\TMright}} (3)
    (3) edge [loop above] node {\TMtrans{\TMstroke}{\TMstroke}{\TMright}} (3)
        edge node {\TMtrans{\TMblank}{\TMstroke}{\TMright}} (4)
    (4) edge [loop below] node {\TMtrans{\TMblank}{\TMstroke}{\TMleft}} (4)
        edge node {\TMtrans{\TMstroke}{\TMstroke}{\TMleft}} (5)
    (5) edge [loop below]  node {\TMtrans{\TMstroke}{\TMstroke}{\TMleft}} (5)
        edge              node {\TMtrans{\TMblank}{\TMblank}{\TMleft}} (6)
    (6) edge [loop below] node {\TMtrans{\TMstroke}{\TMstroke}{\TMleft}} (6)
        edge              node {\TMtrans{\TMblank}{\TMblank}{\TMright}} (1);
\end{tikzpicture}
\]
\caption{يو دوه‌چنده کوونکے ماشين}
\ollabel{fig:doubler}
\end{figure}
\end{ex}

\begin{prob}
يو خپل‌سری ننوت وټاکئ او په~\olref[tur][mac][rep]{ex:doubler} کښې
د دوه‌چنده کوونکي ماشين تشکيلونه پرله‌پسې وڅارئ۔
\end{prob}

\begin{prob}
د~$\{\TMendtape,\TMblank, A, B\}$ الفبا لرونکے ټيورينګ ماشين جوړ
کړئ چې د~$A$ ګانو او~$B$ ګانو هره هغه رشته مني، يعنې پرې درېږي، چې
د~$A$ ګانو شمېر ئې د~$B$ ګانو له شمېر سره برابر وي \emph{او} ټول
$A$ ګانې له ټولو~$B$ ګانو مخکښې راشي؛ او هره هغه رشته ردوي، يعنې
پرې نۀ درېږي، چې د~$A$ ګانو شمېر د~$B$ ګانو له شمېر سره برابر نۀ وي
يا ټولې~$A$ ګانې له ټولو~$B$ ګانو مخکښې رانۀ شي۔ (مثلاً، ماشين بايد
$AABB$ او~$AAABBB$ ومني، خو~$AAB$ او~$AABBAABB$ دواړه رد کړي۔)
\end{prob}

\begin{prob}
د~$\{\TMendtape,\TMblank, A, B\}$ الفبا لرونکے ټيورينګ ماشين جوړ
کړئ چې د~$A$ ګانو او~$B$ ګانو هره رشته~$\alpha$ د ننوت په توګه
واخلي او نقل ئې کړي، څو د~$\alpha\alpha$ په بڼه وت پيدا کړي۔ (مثلاً،
د~$ABBA$ ننوت بايد د~$ABBAABBA$ وت ورکړي۔)
\end{prob}

\begin{prob}
\emph{الفبايي؟}: د~$\{\TMendtape,\TMblank, A, B\}$ الفبا لرونکے
ټيورينګ ماشين جوړ کړئ چې د~$A$ ګانو او~$B$ ګانو يوه متناهي لړۍ د
ننوت په توګه ورکړل شي، نو وګوري چې ټولې~$A$ ګانې د ټولو~$B$ ګانو
کيڼ لوري ته راځي که نه۔ ماشين بايد د ننوت رشته پر پټه پرېږدي، او که
رشته ``الفبايي'' وي ودرېږي، او که نۀ وي تل په کړۍ کښې وچلېږي۔
\end{prob}

\begin{prob}
\emph{الفبايي کوونکے}: د~$\{\TMendtape,\TMblank, A, B\}$ الفبا
لرونکے ټيورينګ ماشين جوړ کړئ چې د~$A$ ګانو او~$B$ ګانو يوه متناهي
لړۍ د ننوت په توګه واخلي او داسې ئې بيا ترتيب کړي چې ټولې~$A$ ګانې
د ټولو~$B$ ګانو کيڼ لوري ته راشي۔ (مثلاً، لړۍ~$BABAA$ بايد په
$AAABB$، او لړۍ~$ABBABB$ بايد په~$AABBBB$ بدله شي۔)
\textbf{د ژباړې د نسخې سمون (OLCMP-042):} سرچينې د ننوت د اخيستلو
او بيا ترتيبولو تر منځ د عطف کلمه پرېښې وه؛ دلته بشپړه جمله راوړل
شوې ده۔
\end{prob}

\end{document}
